\section{The narrow-width comparison}\label{sec:width}

\begin{proposition}\label{prop:narrow}
For every numerical semigroup $S$ of multiplicity $m$, embedding
dimension $e\ge4$, and width at most $m-2$, every field $k$, and every $i\ge0$,
\[
 \beta_i(G_S;k)\le\beta_i(G_{\widetilde S};k).
\]
Both Betti numbers are taken over the polynomial ring on the respective
semigroup's minimal generators.
\end{proposition}

We prove the proposition by covering the possible normalized widths.
Let $S$ have multiplicity $m$, embedding dimension $e\ge4$, and width
$V\le m-2$.  Write its minimal generators as
\[
 S=\langle m,m+da_1,\ldots,m+da_q,m+du\rangle,
 \qquad q=e-2,\quad 0<a_1<\cdots<a_q<u,
\]
where $V=du$, $d\ge1$, $m\ge du+2$, and
$\gcd(a_1,\ldots,a_q,u)=\gcd(m,d)=1$.
This is a change of coordinates for the given generators: their
numerical values and the coefficient field stay fixed.  The comparison
target is still generated by all integers from $m$ through $m+du$.
In particular, any estimate proved at normalized width $u$ must be
transferred to that physical interval.
Use the core, height, cyclic distances, and invariants of
Section~\ref{sec:endpoint}, with $\Gamma=S$.
For brevity put $G=G_S$, $\beta_i=\beta_i(G;k)$,
$Q=Q_C$, $P=P_C$, and $h_n=\dim_k Q_n$.

The endpoint estimates of Section~\ref{sec:endpoint} give $P\le W$,
$R\le\widehat R$, and, writing $b_i=\beta_i(Q;k)$,
\begin{align}
 \beta_i&\le b_i+P\binom q{i-1}+R\binom{q+1}{i-1},\label{eq:width-transfer}\\
 \beta_i+\binom{q+1}{i+1}
 &\le P\binom{q+1}i+\widehat R\binom{q+1}{i-1}.
 \label{eq:width-rows}
\end{align}
The use of the core dimension $P$ in these inequalities includes the
kernel compensation from the endpoint argument; the actual endpoint
quotient can have a different dimension.  Rational coefficients in the
estimates below are weights on dimensions, so the calculations apply
over every characteristic.
The interval completion has, in row $i$, at least
\[
 F_i^-(u)=i\binom u{i+1}
       +\min\left\{\binom u{i-1},\binom ui\right\},
 \qquad 1\le i\le u.
\]
Indeed its physical width is $V\ge u$, and this lower bound increases
with width.  We therefore compare with $F_i^-(u)$, except at the full
arithmetic progression, where a direct comparison is simpler.
Binomial coefficients with an index outside their range are zero.

\subsection{The two smallest widths}

If $u=q+1$, every offset occurs.  The arithmetic progression formula
\cite{APFormula,APCone} bounds every positive row by
$i\binom{u+1}{i+1}$: its correction to $i\binom u{i+1}$ is at most
$i\binom ui$.  For $d\ge2$, $V\ge2u\ge u+1$ supplies this bound.
For $d=1$, $S$ already is its interval completion.

Suppose next that $u=q+2$ and $q\ge3$, and let $h$ be the missing
offset.  The case $q=2$ will be included in the four-generator argument.
For
$B=\{0,\ldots,u\}\setminus\{h\}$, a basis of $Q_n$ is indexed by
\[
 nB\setminus\bigl((n-1)B\cup((n-1)B+u)\bigr).
\]
If $2\le h\le u-2$, then $B+B=[0,2u]$: the missing sums $h,h+u$
are $1+(h-1)$ and $(u-1)+(h+1)$.  Thus $Q$ has Hilbert function
$(1,q,2)$, with quadratic basis at offsets $h,h+u$.  If $h=1$ or
$u-1$, the Hilbert function is $(1,q,1)$.  In every case these basis
elements have height at most one, so $R=0$.

For an interior hole, split the linear variables into blocks of sizes
$r=h-1$ and $s=u-h-1$, so $r+s=q$.  Multiplication pairs each variable
with its unique partner in the same block, producing the corresponding
quadratic basis element.  In the Koszul differential, every nonzero
quadratic row admits a distinct column with its only nonzero entry
$\pm1$ in that row.  Counting the rows whose wedges contain a whole
block, combining this minor with the degree-zero minor, and applying
\eqref{eq:width-transfer} gives
\[
 \beta_i+\delta\binom{q+1}i\le i\binom{q+2}{i+1}+T_i,
 \qquad
 T_i=\binom{s+1}{i-r}+\binom{r+1}{i-s}.
\]
Here $\delta\ge0$ is any integer satisfying $P+\delta\le q+3$;
the same bound allows an empty block at an end hole.
For the following binomial estimates assume an interior hole, so
$r,s\ge1$. End holes are treated separately below.
For $0\le c\le r$, adjoining $r-c$ fixed elements and then enlarging
the ambient set gives
\[
 \binom{s+1}{i-r}\le\binom{q+1}{i-c}.
\]
Apply this with $c=0$ in the first summand of $T_i$, and with $c=1$
in the second.  Pascal gives $T_i\le\binom{q+2}i$.
For $i\ge2$, use $c=2$ in a block of size at least two and $c=1$
in the other block.  Since $r+s\ge3$, this proves
$T_i\le\binom{q+2}{i-1}$; the case $i=1$ is immediate.
Taking $\delta=0$ proves $\beta_i\le F_i^-(u)$ for an interior hole.
For an end hole, without using these binomial estimates, take
$\delta=1$, since $P=q+2$.
The empty-block correction equals $\binom{q+1}i$
plus the two top-row indicators, giving
\[
 \beta_i\le i\binom{q+2}{i+1}
                   +\mathbf1_{i=q}+\mathbf1_{i=q+1}.
\]
Both exceptional rows have target correction at least one.  The minors
use units $\pm1$, so the argument includes characteristic two.

\subsection{Dense widths}

Assume $q\ge4$ and $u=q+1+t$, where $2\le t\le q$.
The prefixes of lengths $2,\ldots,n$ of a surviving word must have
distinct missing residues.  An existing residue would replace a prefix
by a word containing an endpoint; repeated residues would give a
nonempty zero-sum subword and the same contradiction.  Hence
$n-1\le t$ and $H_x\le t$.
The set $A=\{0,a_1,\ldots,a_q\}$ has more than half the residues, so
$A+A$ is the whole group.  Its distance layers consequently give
\begin{equation}\label{eq:width-dense-resources}
 W=q+1+2t,\quad P\le q+1+2t,\quad R\le t,
 \qquad t=2\Longrightarrow R=0.
\end{equation}

Choose one lex quotient with the entire Hilbert function of $Q$.
For this quotient let $S_j(n)$ count degree-$n$ standard monomials in
the first $j$ variables.  Counting minimal generators by their largest
variable gives
$g_j(n)=S_j(n-1)-S_j(n)+S_{j-1}(n)$.
The lex Betti bound, including Pardue's arbitrary-characteristic version
\cite[Theorem~31]{Pardue}, and the stable ideal formula \cite{BHP,EK}, followed
by telescoping in degree, imply
\begin{equation}\label{eq:width-lex}
 b_i\le i\binom{q+1}{i+1}+K\binom{q-1}{i-1},
 \qquad K=\sum_{n\ge2}\sum_{j=1}^qS_{j-1}(n).
\end{equation}
If $\binom{r+1}2\ge2t$ and $r\le q$, all higher standard monomials
are supported on the last $r$ variables.  Indeed each $h_n\le2t$,
and the degree-$n$ monomials in those variables already number at least
$\binom{r+1}2$ and come first among the lex standard monomials.
Count degree-$n+1$ monomials
whose direct divisors are standard, weighting a parent--child edge by
the exponent of the removed variable.  Each child has total weight
$n+1$ and each parent at most $n+r$.
Removing the largest variable from a child gives the prefix-standard
parent count $\sum_jS_j(n)$; its excess over $h_n$ is exactly the
degree-$n$ summand of $K$.  Consequently
\[
 K\le\sum_{n\ge2}\frac{r-1}{n+1}h_n
       \le\frac{2t(r-1)}3.
\]
These weights are rational numbers used to count finite sets.

Put $A_i=\binom{q-1}{i-1}$, $C_i=\binom q{i-1}$,
$B_i=\binom{q+1}{i-1}$.  The first two terms of Vandermonde give
\[
 i\left[\binom{q+t+1}{i+1}-\binom{q+1}{i+1}\right]
 \ge t(q+1)C_i+i\binom t2B_i.
\]
For $t=2$, choose $r=3$, giving $K\le2$; the remaining difference is
at least $(q-3)C_i+iB_i-2A_i\ge0$.
For $t\ge3$, choose $r=t$.  After using $B_i\ge C_i\ge A_i$, the
difference is at least $cC_i$, where
\[
 6c=12(q-4)+5(t-3)^2+(t-3)+6(t-3)(q-t)\ge0.
\]
Equations \eqref{eq:width-transfer}--\eqref{eq:width-lex} therefore
give $\beta_i\le i\binom u{i+1}$ in every positive row.

\subsection{Sparse widths and propagation from the second row}

We first record why one numerical comparison controls all higher rows.
For $e\le u$ define
\[
 U_i(B)=B\binom ei-u\binom{e-1}{i-1}-\binom{e-1}{i+1},
 \quad
 \mathcal T_i(u)=
 \frac{i\binom u{i+1}+u\binom{e-1}{i-1}+\binom{e-1}{i+1}}{\binom ei}.
\]
For $2\le i\le e$, direct binomial identities yield
\[
 \mathcal T_i(e)=e-1+i/e,\qquad
 \mathcal T_i(u+1)-\mathcal T_i(u)
       =i\frac{\binom ui}{\binom ei}+i/e.
\]
The ratio in the increment increases with $i$, since its adjacent ratio
is $(u-i)/(e-i)\ge1$ for $i<e$.  Thus
$\mathcal T_i(u)\ge\mathcal T_2(u)$.
If $P\le B$ and $\widehat R\le B-u$, \eqref{eq:width-rows} bounds
$\beta_i$ by $U_i(B)$.  It follows that
\begin{equation}\label{eq:width-propagation}
 U_2(B)\le2\binom u3
 \quad\Longrightarrow\quad
 \beta_i\le i\binom u{i+1}\quad(2\le i\le e).
\end{equation}

Now assume $e\ge6$ and $u\ge2e-2$.  Set $a_0=e-1$ and
$s=\lceil a_0/2\rceil$.  Kneser's theorem \cite{Kneser} gives
\[
 |nA|\ge\min\{u,a_0+(n-1)s\},\qquad
 W\le B(u):=u+\sum_{j\ge0}(u-a_0-sj)_+.
\]
Indeed a proper stabilizer $H$ meets $A+H$ in $c\ge2$ cosets, with
$c|H|\ge a_0$; its growth contribution $(c-1)|H|$ is at least
$\lceil a_0/2\rceil$.  This also treats composite $u$.
Because $0\in A$, the sets $nA$ increase, and
$nA=\{r:\ell_r\le n\}$.  Summing the missing residues in each layer
gives $W=u+\sum_{n\ge1}(u-|nA|)$ and the displayed bound.
Writing $\Delta=u-a_0=sk+b$, $0\le b<s$, summation gives
\[
 B(u)=u+\frac{\Delta^2}{2s}+\frac\Delta2+\frac{b(s-b)}{2s}
 \le B_*:=u+\frac{\Delta^2}{2s}+\frac\Delta2
                  +\frac{\lfloor s^2/4\rfloor}{2s},
 \qquad B(u)\le\frac{u^2}{e-1}.
\]
The last inequality follows, after clearing denominators, from
\[
 2(u^2-a_0B)=s(2s-a_0)k^2+a_0sk
                  +2(2s-a_0)kb+2b^2\ge0.
\]

The endpoint first-row bounds are
\[
 \beta_1\le\frac{e+3}{3}B-u+\frac{e(e-1)}6,
 \qquad
 \beta_1\le\frac{e+5}{4}B-u+\frac{(e-1)e(e+1)}{24}.
\]
Use the second with $B_*$ for $e=6,7,8,9$.
With $z=u-(2e-2)$, the difference from $\binom u2+1$ is, respectively,
\[
 \frac{z^2+43z+12}{24},\quad 2z+1,\quad
 \frac{3z^2+126z+211}{32},\quad
 \frac{z^2+68z+116}{16}.
\]
For $e\ge10$, the first bound with $B=u^2/(e-1)$ has difference
\[
 \frac{(e-1)(e-10)}2+1+\frac{4e-33}{6}z
                   +\frac{e-9}{6(e-1)}z^2>0.
\]
For row two the same quadratic envelope gives
\[
 2\binom u3-U_2(B(u))\ge
 \frac u6\{2u^2-(3e+6)u+6e-2\}+\binom{e-1}3.
\]
For $e=6$ the right side is $(u-10)(u-3)(u+1)/3\ge0$.
For $e=7+v$, $u=2e-2+z$, the expression in braces is
$2v^2+5vz+12v+2z^2+21z+4>0$.
Apply \eqref{eq:width-propagation} with $B(u)$ to finish every higher row.

\subsection{Five generators}

Here $q=3$.  Widths $u=4,5$ were treated above.
At $u=6,7$, \eqref{eq:width-dense-resources} gives
$(P,\widehat R)\le(8,0),(10,3)$.
The harmonic estimate $b_1\le2P/3+10/3$ gives $\beta_1\le16$ at $u=6$.
At $u=7$, $R=0$ gives $\beta_1\le20$.
If $R>0$, then $H_x\ge3$, so some surviving word has degree at least
four.  For $\tau=\sum_{n\ge3}h_n\ge1$, the degreewise harmonic bound is
\[
 b_1\le6+\frac23(P-4)-\frac16\tau<10.
\]
Integrality gives $b_1\le9$ and $\beta_1\le22$.
For rows $2$--$5$, \eqref{eq:width-rows} gives the sufficient bounds
$(44,31,8,0)$ and $(68,57,22,3)$, below the corresponding capacities
$(46,60,39,11)$ and $(77,126,119,56)$.

For $u\ge8$, Kneser's general increment two gives $4W\le u^2$.
At $u=8$ the unsimplified sum gives $W\le14$ and $U_2(14)=104<112$.
For $u=9+z$, $B=u^2/4$ gives
\[
 2\binom u3-U_2(B)=\frac{2z^3+33z^2+136z+33}{6}>0.
\]
Thus \eqref{eq:width-propagation} treats every row $i\ge2$.

For the first row, an increment smaller than three forces a stabilizer
of order two and exactly two cosets.  Consequently either all proper
increments are at least three, or the original offsets are
\[
 u=2h,\qquad \{a_1,a_2,a_3\}=\{t,h,t+h\},\qquad 0<t<h.
\]
In the first case
$6W\le u^2+u+6$.  The fourth harmonic bound $b_1\le P/2+5$ gives
$\beta_1\le B=(5u^2-7u+90)/12$.
If $A_1=\binom u2+1$, then
$A_1+1-B=((u-8)^2+17(u-8)+6)/12>0$;
integer-valuedness gives $\beta_1\le A_1$.

In the second case denote the three internal variables by $y,z,w$.
The relations $yz=xw$, $z^2=xv$, $zw=vy$ isolate $z$ in $Q$.
For $0\le j\le h$, write $t+j=2K+r$, $r\in\{0,1\}$; then
\[
 y^{h-j}w^j=x^{h-r-K}z^rv^K.
\]
The right side contains an endpoint.  Thus $Q$ is spanned by $z$ and
$y^{n-j}w^j$ with $0\le j\le n<h$, giving
$P\le\binom{h+1}2+1$.
Together with $W\le h^2$, $R\le h^2-2h$, and the third harmonic
bound, this gives
$\beta_1\le B=(11h^2-7h+30)/6$.
For $A_1=\binom{2h}2+1$,
\[
 A_1+1-B=\frac{(h-4)^2+9(h-4)+2}{6}>0.
\]
This proves the required bound for all $h\ge4$ using only a spanning
set of candidate words.

\subsection{Four generators: standard words and higher rows}

Write the two internal offsets as $0<a<b<u$, and their variables as
$y,z$.  Two internal words with the same degree and offset have the
same exponents.  Hence $Q$ is a monomial quotient of $k[y,z]$;
let $D\subset\N^2$ be its finite downset of standard exponents.
Write $\sigma=\dim_k(0:_Q(y,z))$.  Its maximal points form a socle
basis, so its Betti vector is $(1,\sigma+1,\sigma)$.
The endpoint estimates become
\begin{equation}\label{eq:width-four-rows}
 (\beta_1,\beta_2,\beta_3,\beta_4)
 \le(P+\sigma+1+R,\ 2P+\sigma+3R,\ P+3R,\ R).
\end{equation}
Kneser applied to $\{0,a,b\}$ gives $|nA|\ge\min(u,2n+1)$, whence
\[
 P\le W\le1+\lfloor u^2/4\rfloor.
\]
Ordering the $\sigma$ maximal points by increasing first coordinate
forces decreasing second coordinates.  The $j$th point dominates
$(j,\sigma-1-j)$; therefore the whole triangle
$\{(s,t):s+t\le\sigma-1\}$ lies in $D$.
In particular $\binom{\sigma+1}2\le P$ and $\sigma\le u-1$.

A surviving word gives a zero-sum-free residue sequence: a zero-sum
subword can be replaced by endpoint powers.  The long-sequence theorem
of Savchev--Chen \cite{SavchevChen} implies
\begin{equation}\label{eq:width-word-length}
 st>0\Longrightarrow s+t\le\lfloor u/2\rfloor,
 \qquad (s,t)\in D\Longrightarrow s+t\le u-2.
\end{equation}
To see the first assertion, a hypothetical longer mixed word can be
unit-normalized to labels $1,c$, $c\ge2$, with multiplicities $R_0,T_0$
and $R_0+cT_0\le u-1$.  Returning to the original labels gives
$c\alpha=\beta+Ku$, $0\le K\le c-1$, hence
$Y_\alpha^c=x^{c-1-K}Y_\beta v^K$ and $R_0<c$.
Thus $2(R_0+T_0)\le u$, a contradiction.
A pure word of length $u-1$ similarly normalizes its label to one;
the other label gives a forbidden power of degree between $2$ and
$u-1$.  This proves the second assertion.  The unit acts only on this
finite residue sequence.  It follows that $H_x\le u-3$, and hence
$R=0$ for $u\le5$.

From \eqref{eq:width-four-rows} we obtain the rational envelopes
\[
 U_2=5u^2/4-2u+4,\qquad U_3=u^2-3u+4,
 \qquad U_4=u^2/4-u+1.
\]
Put $T_i=i\binom u{i+1}$.  For $u=6+z$,
\begin{equation}\label{eq:width-four-cubic}
 12(T_2-U_2)=4z^3+45z^2+140z+36>0.
\end{equation}
Moreover $U_2-U_3=u^2/4+u$, $U_2-5U_4=3u-1$, and
\[
 8T_3=3(u-3)T_2,\qquad
 10T_4=(u-3)(u-4)T_2.
\]
For $u\ge6$ these imply $U_3\le T_3$ and $U_4\le T_4$.
At $u=5$, the bounds $(P,\sigma,R)\le(7,4,0)$ give higher rows
$(18,7,0)$, below $(25,25,9)$.
At $u=4$, $(P,\sigma,R)\le(5,2,0)$ gives $(12,5,0)$, below
$(12,7,1)$.  Width $u=3$ is a full arithmetic progression.

\subsection{Four generators: the first row}

Three relations in $C$ are
\[
 y^b=x^{b-a}z^a,\qquad
 z^{u-a}=y^{u-b}v^{b-a},\qquad
 y^{u-b}z^a=x^{u-b}v^a.
\]
Set $h=\lfloor u/2\rfloor$ and $\kappa=u-b+a$.
The axes contribute at most $2u-\kappa-1$ standard words.
The mixed words lie in the triangle of size $\binom h2$ in
\eqref{eq:width-word-length}; the third relation removes
\[
 E_h(\kappa)=
 \begin{cases}\binom{h-\kappa+2}2,&\kappa\le h,\\0,&\kappa>h.\end{cases}
\]
Consequently
\[
 P\le2u-\kappa-1+\binom h2-E_h(\kappa)
       \le2u-h-2+\binom h2,\qquad
 \sigma\le\min(h+1,\kappa).
\]
For the last bound, distinct mixed socle points have distinct first
coordinates in $1,\ldots,h-1$, while the socle triangle cannot contain
the forbidden point $(u-b,a)$.
It follows that
\[
 \beta_1\le\binom h2+u+1+\lfloor u^2/4\rfloor.
\]
Subtracting this bound from $\binom u2+1$ gives $h(h-5)/2$ for $u=2h$ and
$(h^2-3h-2)/2$ for $u=2h+1$.
This proves the comparison for $u\ge9$.
For $u=5$, $P\le7$, $\sigma\le3$, $R=0$ give $\beta_1\le11$.

At $u=6$, $H_x>2$ forces a pure word of degree four and offset at
most six.  A mixed word would violate \eqref{eq:width-word-length},
and $z^4$ has offset at least eight.  The remaining word $y^4$ forces
$a=1$; the forbidden power $y^b$ then forces $b=5$.
Thus $(a,b)=(1,5)$; there $yz=xv$ and the pure
fifth-power relations give $(P,\sigma,R)\le(9,2,4)$ and
$\beta_1\le16$.  Otherwise $R=0$ and $\beta_1\le15$.
For $u=7$, the displayed counting bounds give $\beta_1\le22$ unless
$\kappa=4$; the remaining pairs are $(1,4),(2,5),(3,6)$.
They satisfy, respectively, $z^2=yv$, $yz=xv$, $y^2=xz$, which give
$(P,\sigma)\le(7,2),(9,2),(7,2)$.
With $R\le6$ these also suffice.
For $u=8$ only $\kappa=4,5,6$ remain after counting.  The necessary
strip bounds are
\[
\begin{array}{c|c|c}
 \kappa&(a,b)&(P,\sigma)\text{ at most}\\\hline
 4&(1,5),(3,7)&(14,4)\\
 5&(1,4),(2,5),(3,6),(4,7)&(9,2)\\
 6&(1,3),(5,7)&(14,3)\\
 6&(3,5)&(9,2).
\end{array}
\]
For $\kappa=4$, the axes have at most nine points and the six-point
mixed triangle loses the forbidden point $(3,1)$ or $(1,3)$.
For $\kappa=5$, one square has an endpoint factor and the other pure
power has degree at most five; the axes and mixed words contribute at
most six and three points.  For $\kappa=6$, the outer pairs have a
forbidden cube and a forbidden seventh power, leaving at most nine
axis points and five mixed points.  In the middle pair $yz=xv$ removes
all mixed words.  The corresponding strip widths bound the number of
socle corners by two or three, as recorded in the table.
These give $\beta_1\le29$ using $R\le9$.
The lists follow by solving $b-a=u-\kappa$ and imposing
$\gcd(a,b,u)=1$.

At $u=4$, $d\ge2$ is immediate from the coarse bound $\beta_1\le8$
and $V\ge8$.  For $(a,b)=(1,2),(2,3)$ both squares vanish in $Q$.
The resulting two-by-two downset satisfies $P+\sigma\le5$, so
$\beta_1\le6$.  It remains to consider $d=1$, $(a,b)=(1,3)$.

\subsection{The remaining width-four family}

Here $m\ge6$.  Use the four homogeneous relations
\begin{equation}\label{eq:width-four-core}
 y^3-x^2z,\quad yz-xv,\quad xz^2-y^2v,\quad z^3-yv^2.
\end{equation}
Reduction decreases the internal weight $\deg_y+2\deg_z$ and leaves
the columns $x^Lv^Kb$ with $b\in\{1,y,y^2,z\}$, and $z^2v^K$.
Their degree--offset pairs are distinct, so they are independent in $C$.

Let an anchor $b$ have degree $c$ and offset $n$.
A competing free column with degree larger by $j\ge1$ has residue
$r_j=(n-jm)\bmod4$ and coordinates
\[
 K'=K-\kappa_j,\qquad L'=L-\Delta_j,\quad
 \kappa_j=\left\lceil\frac{jm-n}{4}\right\rceil,\quad
 \Delta_j=\lambda(r_j)-c-\kappa_j-j,
\]
where $(\lambda(0),\lambda(1),\lambda(2),\lambda(3))=(0,1,2,1)$.
Since $4c-n\ge0$ and $4\lambda(r)-r\le6$,
$4\Delta_j\le6-j(m+4)<0$.
Thus competition requires exactly $K\ge\kappa_j$.
A restricted column has offset equal to four times its degree minus
two, and cannot be a longer competitor when $m\ge6$.
Since $\kappa_j$ increases with $j$, the full condition reduces to
\[
 K\ge\theta_n,\qquad
 \theta_n=\max\left\{\left\lceil\frac{m-n}{4}\right\rceil,0\right\}.
\]
The defining ideal of $G$ is therefore generated by
\eqref{eq:width-four-core} and the five tails
$bv^{\theta_n}$ for $b=1,y,y^2,z,z^2$.
Indeed reduction spans the quotient by the surviving columns, whose
images are independent semigroup monomials in their exact order.

For $m=6$, two core relations and four tails suffice, giving
$\beta_1\le6<9$.
For $m>6$, write $m=4Q+\rho$.  In anchor order $1,y,y^2,z,z^2$,
the five thresholds are
\[
\begin{array}{c|ccccc}
 \rho&\theta_0&\theta_1&\theta_2&\theta_3&\theta_6\\\hline
 0&Q&Q&Q&Q&Q-1\\
 1&Q+1&Q&Q&Q&Q-1\\
 2&Q+1&Q+1&Q&Q&Q-1\\
 3&Q+1&Q+1&Q+1&Q&Q.
\end{array}
\]
A tail is redundant whenever a divisor anchor has a threshold no
larger than its own.  Retain all four core relations and, according to
$\rho$, respectively two, four, four, or two tails.  Hence
\[
\begin{array}{c|rrrr}
 \rho&0&1&2&3\\\hline
 \text{number of generators at most}&6&8&8&6\\
 F_1(m,4)&7&10&9&8.
\end{array}
\]
This completes the first row for every $m$ in the remaining family.

The preceding cases cover $e=4$ and all $u\ge3$.
For $e\ge6$, the ranges $u=e-1,e$, $e+1\le u\le2e-3$, and
$u\ge2e-2$ are covered by the smallest, dense, and sparse width
arguments; the five-generator discussion covers $e=5$.
Finally $\beta_0=1$ on both sides and $\beta_i(G)=0$ for $i>e$.
Thus every narrow source satisfies the Betti comparison with its
literal interval completion, over the original field.
